% ============================================================
% 01_introduction.tex  (revised 2026-10)
% ============================================================

Global energy systems are being reshaped by the diffusion of solar photovoltaic (PV) generation.
Cumulative PV capacity reached approximately $1{,}865~\mathrm{GW}$ by the end of~2024, with a
record $452~\mathrm{GW}$ added during that year~\cite{irena2025}. In large interconnected grids,
forecast errors are absorbed by reserves and markets; in small and off-grid microgrids, where PV
competes with a few thermal generators, they translate directly into spinning reserve, fuel
consumption and curtailment. Forecasts covering the next hours to the next day are therefore an
operational input of microgrid dispatch.

Machine-learning forecasters and their ensembles are routinely reported to outperform single
models, but three aspects of their evaluation are often fragile. First, the combination strategy
or meta-learner is sometimes fitted on the period used for the final scores, which inflates the
reported gains~\cite{kaufman2012leakage,kapoor2023leakage}. Second, the measured PV output is
taken as the physical resource although it may include operator decisions such as curtailment.
Third, gains are frequently reported on normalised targets, a single seed and a single site,
without tests that account for the strong dependence between overlapping multi-step forecasts.

This paper revisits the Adaptive Honest Selection Ensemble (AHSE), first proposed for the Yulara
plant in Australia, with these issues in mind. We found that the Yulara series used in earlier
studies, including the first version of this work, is the output of a site that is curtailed
around noon; the corresponding claims are withdrawn and the study is rebuilt on uncurtailed data.
The contributions are:

\begin{enumerate}
  \item \textbf{A curtailment-aware benchmark.} We document the curtailment of the Desert Gardens
  site from the public per-site data (about half of its available energy), and use the four
  uncurtailed sites of the same mini-grid as target, with a second, grid-connected site in a
  humid subtropical climate (NIST, USA) obtained from a public data lake.

  \item \textbf{Honest per-horizon-block selection.} AHSE selects, for each block of lead times,
  a single model or an average of the best models, and accepts a combination only when the lower
  bound of a day-level bootstrap interval of its validation gain is positive. All decisions use
  the validation period only.

  \item \textbf{Leak-free weather forecasts.} Archived GFS forecasts enter the tree members and
  the candidate pool using only the run published before the issue time, with a conservative
  five-hour delay.

  \item \textbf{Calibrated uncertainty.} Out-of-fold, Mondrian and adaptive conformal intervals
  are calibrated without touching the test period, optionally by a weather-forecast variable.

  \item \textbf{Transparent evaluation.} Errors in power and as a share of capacity, a
  squared-error skill score against day-ahead persistence, five seeds, day-level bootstrap and
  Diebold--Mariano tests, and a sensitivity analysis of the validation design; all scripts and
  result tables are public.
\end{enumerate}

Section~\ref{sec:related} reviews related work. Section~\ref{sec:methodology} defines the
forecasting task, the selection protocol, the use of weather forecasts and the prediction
intervals. Section~\ref{sec:setup} describes the sites, the curtailment analysis and the
evaluation protocol. Section~\ref{sec:results} reports and discusses the results and
Section~\ref{sec:conclusion} concludes.
